\documentclass{jsarticle}
\usepackage[dvipdfmx,final]{graphicx}
\usepackage{chemfig}
\usepackage{amsmath}
\usepackage{amssymb}
\begin{document}
\title{}
\author{}
\date{}
\maketitle


           Amulozipin is destreamed with entrance of 
	acasic record in splitual mind to respect 
	from surface of heart on metal gardient,
	this gard of heart destgusted from this medicine and
	desease with posion of blood without pressure of pulse,
	this despoit of magnity of field on reverse of person with
	healty of body and splitual cells. Moreover this medicine
	of semilar with being kept with amel of fullnitrazepame
	stimulate with selen with being haven to cover with Oxiden
	into gus from solid that being in kalumm and calsium to
	emerge with magnity convert into being deceived with blood
	of pressure with essence of type of reverse's persons.



        アムロジピンは、血流を停滞させる働きがあり、ほっとする効果が
    精神に作用する。さらに、心臓の鉄分を取り除くと、磁場も発生せずに、
    普通は、鉄分を取ると、心筋梗塞や付随的な作用をするが、
    この磁場を取り除くと、真逆の人は、健康的な精神を保つことができる。
    脳にも同じ作用を磁場によって精神にも作用して、アメルのフルニトラゼパム
    は、セレンは固体だが、酸素を付加すると気体になり、さらに、このセレンに
    カリウムとカルシウムを磁場発生に使うと、真逆は磁場の発生を消すことが
    できて、健康的な精神になる。



        塩分とカルシウムを心臓にも、同じ原理が働き、塩分とカルシウムのバランス
    によって、アムロジピンと同じ作用をする。ドーパミンブロックは、
    このアムロジピンと合わすのに、西洋の漢方薬と同じく、元型と中庸、
    定型の原理が大切らしく、脇本先生は、漢方薬の不思議さを西洋医学で
    見えているらしい。荘六先生が、西洋的な漢方薬の処方に長けている。
    ２００３年に葉子先生がリスペリドンを２ｍｇ出している人に、
    荘六先生は、１ｍｇで、その人にリスペリドンの事情を言っているのを、
    私は、今年どういうことか知らされた。１ｍｇは相当安全らしい。
    芳香族化合物で、セレンやマグネシウムの希釈を二量体で実現している上に、
    その金属の作用を高める働きを、この二量体はしている。
    典型的な有機化合物では、お酢の酢酸が、この二量体を実現している。
    ナトリウムは、程々の量が必要なために、私が守れないために、先生は、
    大変ならしい。取らな過ぎも大変らしく、運動すると、私の体内のナトリウム
    の量が変わるので、アムロジピンは、もっと増えてもいいらしい。女医に
    教えられている。他者の思考をわかる原理は、体温低温と血圧低圧であり、
    上がらない思考で見るのが安全らしい。広島大学医学部から、入院時から、
    母以外のマクロビが始まっていた。仁風荘でもされている。蔵王病院でも
    されている。
    磁場のことは、ネコにノミがつかなくなっていることからわかった。



        叡智を取ると叡智が消えるらしい。これを先生たちが言えないから、
    エチゾラムとブロチゾラムで言いたかったらしい。
    ナトリウムを減らすとは、真逆は相当な効果を、できたら期待できるらしい。
    この理論は、スコトーマだらけであり、心臓に鉄分がどう作用しているかを
    知り、わかったことでもある。
    荘六先生は、私にわかりやすく、ナトリウムを減らすのを防ぐカルシウム
    拮抗薬を処方している。



        夢で、睦世お姉さんのフルニトラゼパムが、磁性体の真逆の作用と、
    アカシックレコードが消えること、二回目が、ナトリウムが心臓に
    どう作用して、鉄分を取り除くと、磁場が消えて、健康的な精神と身体へと
    作用をするのかが、アムロジピンとわかったことである。
    普通は、鉄分は健康に必須であり、鉄分は妄想を消すのに使われるが、
    低体温と低血圧が実現したら、磁場を消しても大丈夫である。
    普通は、磁場は身体の調子を整えるのに使われるが、
    真逆は、ナトリウムと同様に消しても健康には差し支えない。
    逆に、磁場がないとどう心臓に鉄分が作用するかを調べると、
    この部分の機知がわかると、ノーベル賞に匹敵するのは、
    ナトリウムとカリウム、カルシウム、そして、セレンの精神と身体に
    対する作用がわかると、納得される理論とできる。
    心臓にナトリウムが多めに作用すると、痙攣らしく、
    そのために、アムロジピンは、鉄分を心臓に少ししか送らなく、
    そのために、血流が停滞して、磁場が体全体になり、電磁場が
    体から流れなくなり、磁場が消える。電磁場は、電場と磁場の
    偏光体であり、片方がなくなると、電磁場として、存在できない。
    精神は、化学反応の相互作用と脇本先生に教えられている。
    精神の単体は、軸索を飛び飛びに伝わるミンリエ輪根での
    電磁気の流れの結果であり、
    血管の微小管の量子振動による現象の賜物でもあり、
    磁場を消すと、ずっと瞑想が達成できることでもあり、
    神経の高まりを鎮めるのに、磁場を消すと、精神と身体が真逆で健康
    になることが、この理論の結論でもある。


  $$ \chemfig{*6(=-(-[,1.2]C?(=[:-90,1.2]O))=(-[,1.2]C(=[:90,1.2]O)
   -[:-40,1.6]O?)-=-)} $$
  $$  \chemfig{*6(=(-[:-90,0.6]OCa)-=-=-[:0,0.6]-)}
   \chemfig{(-[:0,0.6]N=N(-[:-90,0.6]Cl))}
   \chemfig{(-[:180,0.6])*6(=(-[:-90,0.6]OH)-=-=(-[:120,0.6]OH)-)}\chemfig{
	   *6(-O---(-[:90,0.6]R2)--)}\chemfig{(-[:-30,0.6]*6(-=-(-[:-30,0.6]OH)
   =(-[:30,0.6]OK)-(-[:90,0.6]R1)=-))} $$
 
\end{document}
